\section{引言}
\label{sec:intro}
% 五步，见 knowledge/paper-anatomy.md 第 2 节。不设小标题，用 \paragraph{}。

% (1) 代价。一个真实事故或调查数字，再收成开放问题。
<问题>普遍存在且代价高昂~\cite{placeholder2024}。<年份>，<带数字的具体事故>~\cite{placeholder2024}。如何在只有<约束>时<目标>，仍是一个开放问题。

% (2) 视角转换与命名的性质。
然而<传统目标>并不能保证<下游目标>。在这一视角下，两个性质共同决定<结果>。\emph{<性质 A>}度量<……>。\emph{<性质 B>}度量<……>。图~\ref{fig:motivation}与例~\ref{ex:motivation}具体展示了二者的取舍。

\begin{figure}[t]
\centering
\fbox{\parbox{0.9\linewidth}{\centering\vspace{1.2cm}动机图，见 knowledge/figure-archetypes.md 第 1 节。\vspace{1.2cm}}}
\caption{<在什么条件下比较了什么，每个面板一句自然的话>。}
\label{fig:motivation}
\end{figure}

% (3) 真实数据集上的例子与 Observation。
\begin{example}[<短标题>]
\label{ex:motivation}
我们以公开数据集<数据集>~\cite{placeholder2024}的一个子集为例。<每行是什么、任务和模型>。<划分与规模>。<搜索空间以及图~\ref{fig:motivation}展示了什么>。
\end{example}

\textbf{观察。}由这个例子可以得到两点观察。
\textbf{(1) <一句话论断>。}<带数字的证据>。
\textbf{(2) <一句话论断>。}<带数字的证据>。

% (4) 挑战。一个承上启下的问句，再列编号的挑战。
\paragraph{挑战。}上述观察说明<要求>。那么现有方法为何没有解决<问题>？原因可归结为三项挑战。
\textbf{(1) <收益依赖任务>。}<原因，结合例~\ref{ex:motivation}>。现有的<方法族>~\cite{placeholder2024}缺少<能力>。
\textbf{(2) <决策空间组合爆炸>。}<原因>。
\textbf{(3) <决策难以迁移>。}<原因>。

% (5) 贡献。第一条承载核心论点，最后一条讲实验。
\paragraph{贡献。}针对上述不足，本文提出 \sys{}，<一句话定义>，贡献如下。
(1) \textbf{<名称>。}<做了什么、为什么有效，定理~\ref{thm:guarantee}>。
(2) \textbf{<名称>。}<做了什么、为什么有效，第~\ref{sec:method}~节>。
(3) \textbf{实验结果。}我们在 \num{datasets} 个数据集上与 \num{baselines} 个基线比较。\sys{} <头条结果>，代价为<代价>。
